% !TeX program = xelatex
% 直接编译本文件即可查看常用写法；替换题面和题解开始自己的作业。
\documentclass[nowatermark]{homework}
\newcommand{\hwname}{XXX}
\newcommand{\hwid}{在此填写学号}
\newcommand{\hwemail}{}
\newcommand{\hwclass}{示例课程}
\newcommand{\hwtype}{作业}
\newcommand{\hwnum}{1}
\begin{document}
\maketitle
本文件按功能展示模板用法。每项先给出排版效果，再给出对应写法。
\newcommand{\demoheading}[1]{\par\bigskip\noindent{\large\bfseries\heiti #1}\par\medskip}
\newcommand{\demolabel}[1]{\par\medskip\noindent{\bfseries\heiti #1}\par\smallskip}
\demoheading{1. 题目、题面与题解}

\demolabel{效果}
\question
\begin{problem}
在这里填写题面。
\end{problem}
\solution
在这里填写题解。
\demolabel{写法}
\begin{lstlisting}[language={[LaTeX]TeX},basicstyle=\ttfamily\footnotesize,numbers=none]
\question
\begin{problem}
在这里填写题面。
\end{problem}
\solution
在这里填写题解。
\end{lstlisting}
\newpage
\demoheading{2. 字母小问：alphaparts}
同一题内字母编号连续。题解需要重新从 (a) 开始时，先重置 \texttt{partCounter}。
\demolabel{效果}
\setcounter{partCounter}{0}
\begin{alphaparts}
  \questionpart 第一小问。
  \questionpart 第二小问。
\end{alphaparts}
\demolabel{写法}
\begin{lstlisting}[language={[LaTeX]TeX},basicstyle=\ttfamily\footnotesize,numbers=none]
\setcounter{partCounter}{0}
\begin{alphaparts}
  \questionpart 第一小问。
  \questionpart 第二小问。
\end{alphaparts}
\end{lstlisting}
\demoheading{3. 独立分点：plainparts}
每次从 1 开始，不带题号；可以放在题解中，也可以嵌入字母小问。
\demolabel{效果}
\begin{plainparts}
  \plainpart 整理输入。
  \plainpart 展开推导。
  \plainpart 写出结论。
\end{plainparts}
\demolabel{写法}
\begin{lstlisting}[language={[LaTeX]TeX},basicstyle=\ttfamily\footnotesize,numbers=none]
\begin{plainparts}
  \plainpart 整理输入。
  \plainpart 展开推导。
  \plainpart 写出结论。
\end{plainparts}
\end{lstlisting}
\newpage
\demoheading{4. 普通列表：enumerate}
普通列表使用标准的 \texttt{\textbackslash item}，不需要设置编号格式。
\demolabel{效果}
\begin{enumerate}
  \item 第一步。
  \item 第二步。
\end{enumerate}
\demolabel{写法}
\begin{lstlisting}[language={[LaTeX]TeX},basicstyle=\ttfamily\footnotesize,numbers=none]
\begin{enumerate}
  \item 第一步。
  \item 第二步。
\end{enumerate}
\end{lstlisting}
\demoheading{5. 在小问中分点作答}

\demolabel{效果}
\setcounter{partCounter}{0}
\begin{alphaparts}
  \questionpart 第一小问的答案：
    \begin{plainparts}
      \plainpart 第一步。
      \plainpart 第二步。
    \end{plainparts}
  \questionpart 第二小问的答案。
\end{alphaparts}
\demolabel{写法}
\begin{lstlisting}[language={[LaTeX]TeX},basicstyle=\ttfamily\footnotesize,numbers=none]
\setcounter{partCounter}{0}
\begin{alphaparts}
  \questionpart 第一小问的答案：
    \begin{plainparts}
      \plainpart 第一步。
      \plainpart 第二步。
    \end{plainparts}
  \questionpart 第二小问的答案。
\end{alphaparts}
\end{lstlisting}
\newpage
\demoheading{6. 多级小问：subparts}
下级编号继承当前小问，支持 a.1、a.2、a.1.1。包括最外层在内，最多四层。
\demolabel{效果}
\setcounter{partCounter}{0}
\begin{alphaparts}
  \questionpart 第一小问。
    \begin{subparts}
      \subpart 第一个下级问题。\label{part:child}
        \begin{subparts}
          \subpart 更深一层。
        \end{subparts}
      \subpart 第二个下级问题。
    \end{subparts}
  \questionpart 第二小问。
    \begin{subparts}
      \subpart 重新从 b.1 开始。
    \end{subparts}
\end{alphaparts}
引用下级小问：\ref{part:child}。
\demolabel{写法}
\begin{lstlisting}[language={[LaTeX]TeX},basicstyle=\ttfamily\footnotesize,numbers=none]
\setcounter{partCounter}{0}
\begin{alphaparts}
  \questionpart 第一小问。
    \begin{subparts}
      \subpart 第一个下级问题。\label{part:child}
        \begin{subparts}
          \subpart 更深一层。
        \end{subparts}
      \subpart 第二个下级问题。
    \end{subparts}
  \questionpart 第二小问。
    \begin{subparts}
      \subpart 重新从 b.1 开始。
    \end{subparts}
\end{alphaparts}
引用下级小问：\ref{part:child}。
\end{lstlisting}
\newpage
\demoheading{7. 题面里的三线表：threetab}
列格式不加竖线，表头后用 \texttt{\textbackslash midrule}；顶线和底线由环境生成。
\demolabel{效果}
\begin{problem}
下表列出输入数据。
\begin{threetab}{lr}
项目 & 数量 \\
\midrule
甲 & 10 \\
乙 & 20 \\
\end{threetab}
\end{problem}
\demolabel{写法}
\begin{lstlisting}[language={[LaTeX]TeX},basicstyle=\ttfamily\footnotesize,numbers=none]
\begin{problem}
下表列出输入数据。
\begin{threetab}{lr}
项目 & 数量 \\
\midrule
甲 & 10 \\
乙 & 20 \\
\end{threetab}
\end{problem}
\end{lstlisting}
\newpage
\demoheading{8. 数学归纳证明：induction}

\demolabel{效果}
\begin{induction}
  \basecase 当 $n=1$ 时，等式成立。
  \indhyp 假设 $n=k$ 时等式成立。
  \indstep 利用假设推导 $n=k+1$ 的情形。
\end{induction}
\demolabel{写法}
\begin{lstlisting}[language={[LaTeX]TeX},basicstyle=\ttfamily\footnotesize,numbers=none]
\begin{induction}
  \basecase 当 $n=1$ 时，等式成立。
  \indhyp 假设 $n=k$ 时等式成立。
  \indstep 利用假设推导 $n=k+1$ 的情形。
\end{induction}
\end{lstlisting}

% 有引用时，在正文使用 \citep 或 \citet，再启用文末参考文献。
% \bibliography{ref}
\end{document}
